% testing_deep.tex — Deep-dive on progressive testing methodology
% Appendix L

\section{Progressive Testing Methodology}\label{sec:testing-deep}

This appendix documents the five-tier verification framework that governed all testing throughout the project, from desktop simulation through to autonomous flight.  The framework is motivated by a fundamental asymmetry in autonomous UAV development: a logic error that manifests as a misplaced pixel in simulation can, on real hardware, command a 2\,kg aircraft into an unrecoverable trajectory.  The cost of discovering a defect rises steeply with integration level---a sign-convention bug caught in simulation costs minutes; the same bug discovered during autonomous flight costs days of rework, potential hardware damage, and a lost flight slot that a five-person university team cannot easily reschedule.  The following subsections detail the philosophy, tier definitions, test script organisation, defect discovery data, V-model mapping, and cost--risk gradient that together form the project's verification and validation strategy.

\subsection{Philosophy: Never Skip a Tier}\label{sec:test-philosophy}

The governing principle is borrowed from safety-critical avionics: \emph{each verification tier gates the next, and no gate may be bypassed}. A failure at Tier~2 (bench) blocks all flight testing until the root cause is resolved, regardless of schedule pressure. This discipline is counterintuitive on a time-constrained student project, where the temptation to ``just try it in the air'' is strongest precisely when time is shortest. The justification is empirical: of the 15 defects discovered during this project's verification campaign, 12 were caught at Tiers~1--3 (Table~\ref{tab:bug-cost-appendix}), where the turnaround time from discovery to fix was under one hour. Had those defects propagated to flight, each would have consumed an entire field session.

The framework also enforces a separation of concerns. Each tier tests a \emph{specific integration boundary}: Tier~1 tests software logic; Tier~2 tests the software--hardware interface; Tier~3 tests the hardware--environment interface; Tier~4 tests the perception pipeline under operational conditions; Tier~5 tests closed-loop autonomy. This decomposition ensures that when a failure occurs, its root cause is immediately localisable to the boundary most recently crossed.

% ─────────────────────────────────────────────────────────────────────
\subsection{The Five Tiers in Detail}\label{sec:tiers-detail}

\begin{figure}[H]
\centering
\small
\setlength{\tabcolsep}{2pt}
\begin{tabular}{@{}c@{\hspace{6pt}}c@{\hspace{6pt}}c@{\hspace{6pt}}c@{\hspace{6pt}}c@{}}
\fbox{\parbox{2.4cm}{\centering\textbf{Tier 1}\\SITL\\Simulation}} &
$\xrightarrow{\text{gate}}$ &
\fbox{\parbox{2.4cm}{\centering\textbf{Tier 2}\\Bench Test\\(no props)}} &
$\xrightarrow{\text{gate}}$ &
\fbox{\parbox{2.4cm}{\centering\textbf{Tier 3}\\Passive\\Flight}} \\[8pt]
& & & & $\downarrow$ \text{gate} \\[4pt]
& & \fbox{\parbox{2.4cm}{\centering\textbf{Tier 5}\\Full\\Autonomous}} &
$\xleftarrow{\text{gate}}$ &
\fbox{\parbox{2.4cm}{\centering\textbf{Tier 4}\\Autonomous\\Log-Only}}
\end{tabular}
\caption{Tier progression. Each arrow represents a gate: the downstream tier is blocked until all upstream tests pass.}\label{fig:tier-flow}
\end{figure}

\subsubsection{Tier 1: SITL Simulation (Zero Hardware Risk)}

ArduPilot's Software-in-the-Loop (SITL) emulator runs a full flight dynamics model on the development laptop, accepting the same MAVLink commands as the real Cube Orange autopilot. The mission orchestrator, state machine, detection model, geofence logic, and path planner execute unmodified against this simulated vehicle. A synthetic camera view is generated by extracting a field-of-view-correct crop from a georeferenced satellite image (\texttt{map.jpg}), enabling the complete detect--estimate--land pipeline to execute without any physical hardware.

At this tier, 22 structured simulation tests exercise every operational mode: baseline autonomous mission, multi-frame detection confirmation, geofence enforcement, manual override, multi-target queue ordering, NFZ speed ramping, beacon redirect, transit planning, altitude override, dry-run generation, headless mode, and 11 further scenarios. The interactive simulator (\texttt{simple\_simulator.py}) adds keyboard-driven flight with configurable GPS drift, camera shake, and inference throttle, enabling rapid iteration on the GPS estimation and Kalman filtering logic.

A separate \texttt{-{}-dry-run} mode walks the state machine through all transitions without connecting to SITL, printing the full waypoint set and generating a pattern visualisation. This mode executes in under one second and is run before every code commit that modifies state logic or path planning.

\paragraph{Defects caught.} State machine logic errors (incorrect transition guards, missing timeout handlers), geofence sign-convention bug (\texttt{cv2.pointPolygonTest} returns positive for \emph{inside}, opposite to the initial assumption), repulsive-force direction inversion at NFZ boundary, barometric drift causing infinite landing loops (resolved with a 90\,s absolute timeout), and duplicate-detection queueing of the same target from consecutive frames (resolved with a 5\,m reject radius).

\subsubsection{Tier 2: Bench Test (Pi + Cube + Camera, No Props)}

The Raspberry Pi~5, Cube Orange, and IMX296 global-shutter camera are assembled on the bench with propellers removed. Three numbered scripts test with increasing integration:

\begin{enumerate}
    \item[\texttt{0a}] \textbf{Cube commands} --- sends individual MAVLink commands (mode switch, parameter read, arm pre-check) and verifies each \texttt{COMMAND\_ACK}. Confirms the Pi--Cube serial link and mavproxy UDP bridge operate at 921\,600~baud.
    \item[\texttt{0b}] \textbf{Bench mission} --- executes the full command sequence (arm, takeoff, waypoint navigation, land) and logs every ACK. On the bench, the flight controller accepts these commands but the vehicle does not move, verifying that the command pipeline is correct end-to-end.
    \item[\texttt{0c}] \textbf{Feedback test} --- runs the vision-to-GPS pipeline while the operator walks the drone past a printed dummy. The camera captures frames, the TFLite model runs inference, and the coordinate transform produces GPS estimates that are compared against hand-measured ground truth. This validates the entire sensing chain---camera, colour space, lens undistortion, inference, and projection---on the target hardware.
\end{enumerate}

Additional bench scripts include a 50-run inference benchmark (\texttt{benchmark.py}), a comprehensive multi-model benchmark with system profiling (\texttt{benchmark\_full.py}), GPS satellite tracking (\texttt{gps\_test.py}), and a lens distortion calibration using a checkerboard target (\texttt{lens\_calibrate.py}).

\paragraph{Defects caught.} BGR/RGB colour-space mismatch (the IMX296 outputs BGR data despite picamera2's \texttt{RGB888} label, causing inverted colour channels that degraded detection confidence from 0.97 to 0.3), FOV miscalibration (focal length corrected from 7.0\,mm to 5.46\,mm after ruler-based measurement), Python~3.13 compatibility breakage (pyserial's serial read is broken on 3.13, requiring a mavproxy UDP bridge as a workaround; \texttt{tflite-runtime} package replaced by \texttt{ai-edge-litert}), and camera-in-use resource conflicts when multiple scripts attempt concurrent access.

\subsubsection{Tier 3: Passive Flight (Pilot Flies, Pi Watches, Zero Commands)}

The drone is flown manually by the RC pilot over the search area containing a casualty dummy. The Pi runs the full vision pipeline---camera capture at $1456 \times 1088$, YOLOv8n inference via TFLite, lens undistortion, GPS coordinate estimation---and streams the annotated video to a browser-based ground station. Critically, the Pi sends \emph{zero flight commands} to the Cube. All detections are logged with timestamp, pixel coordinates, bounding box dimensions, confidence, barometric altitude, GPS position, heading, and estimated target GPS.

Post-flight analysis of these logs establishes:
\begin{itemize}
    \item \textbf{Detection altitude ceiling} --- the maximum altitude at which the model reliably detects the dummy (target: 30\,m, to be confirmed).
    \item \textbf{False-positive rate} --- detections triggered by terrain features, shadows, or equipment under real lighting and motion conditions.
    \item \textbf{GPS estimation accuracy} --- scatter of estimated target positions versus surveyed ground truth, quantified as CEP50 and CEP95.
    \item \textbf{Inference latency under load} --- confirmed at 208\,ms (4.8~FPS) on Pi~5 with XNNPACK, including 1.5\,ms for lens undistortion.
\end{itemize}

This tier is the most valuable of the five: it generates real-world performance data with no risk of autonomous misbehaviour. The six experiment scripts in \texttt{tests/day\_1\_experiments/} (Section~\ref{sec:experiment-scripts}) are designed specifically for this tier, each producing a CSV file suitable for direct statistical analysis.

\subsubsection{Tier 4: Autonomous Search with Log-Only Detection}

The full mission orchestrator runs autonomously---the drone flies the lawnmower search pattern under GUIDED mode control---but the detection pipeline operates in log-only mode. When the model identifies a candidate, the system records the detection and annotates the ground station display but does \emph{not} transition to the CENTERING or DESCENDING states. The drone completes the search pattern and returns to launch.

Post-flight review of the detection log answers the operational question: \emph{would the system have found the target, and would it have acted on the correct detection?} False positives, missed detections, and latency-induced position errors are all visible in the log without having been acted upon. Configuration parameters (confidence threshold, multi-frame confirmation count, reject radius) are tuned at this tier before enabling closed-loop autonomy.

\subsubsection{Tier 5: Full Autonomous Mission}

All systems are enabled. The drone executes the complete state machine: INIT, CONNECTING, ARMING, TAKEOFF, SEARCH (lawnmower pattern), CENTERING (visual servo onto target), VERIFY (operator confirms or rejects via ground station), APPROACH (offset landing approach), LANDING, and DONE. A DESCENDING state for altitude reduction before VERIFY exists in the code but is currently bypassed; the drone transitions directly from CENTERING to VERIFY at search altitude. The operator retains veto authority at the VERIFY gate and the RC kill switch provides a hardware-level override independent of all software.

% ─────────────────────────────────────────────────────────────────────
\subsection{Test Script Organisation}\label{sec:test-org-detail}

Table~\ref{tab:test-categories} maps the six test script directories to the verification tiers they support, with representative scripts and total counts.

\begin{table}[H]
\centering
\caption{Test script categories and their tier coverage.}
\label{tab:test-categories}
\small
\begin{tabular}{@{}llcl@{}}
\toprule
\textbf{Category} & \textbf{Directory} & \textbf{Scripts} & \textbf{Tiers Served} \\
\midrule
Hardware     & \texttt{tests/hardware/}           & 8  & 2 (bench) \\
Flight       & \texttt{tests/flight/}             & 8  & 2, 3, 4, 5 \\
Diagnostics  & \texttt{tests/diagnostics/}        & 5  & 2, 3, 4, 5 \\
Calibration  & \texttt{tests/calibration/}        & 5  & 2, 3 \\
Experiments  & \texttt{tests/day\_1\_experiments/} & 6  & 3, 4 \\
Laptop       & \texttt{tests/laptop/}             & 16 & 1 (simulation) \\
Unit         & \texttt{tests/unit/}               & 6  & 1 \\
Automated    & \texttt{tests/automated/}          & 4  & 1 \\
\midrule
\multicolumn{2}{@{}l}{\textbf{Total}} & \textbf{58} & \\
\bottomrule
\end{tabular}
\end{table}

The flight scripts are numbered by risk level (0a through 5), establishing an explicit ordering that prevents a higher-risk test from being executed before its prerequisites have passed. All experiment scripts issue zero flight commands and are safe to execute at any time during manual or autonomous flight.

% ─────────────────────────────────────────────────────────────────────
\subsection{Experiment Scripts for Quantitative Validation}\label{sec:experiment-scripts}

Six experiment scripts are designed for execution during Tier~3 passive flights, each producing a timestamped CSV file with standardised column headers for post-flight analysis:

\begin{enumerate}
    \item \textbf{Altitude sweep} (\texttt{altitude\_sweep.py}) --- the pilot hovers above the dummy at 10, 15, 20, 25, and 30\,m for 30\,s each. The script bins every frame by barometric altitude into 5\,m buckets and computes detection rate, mean confidence, and expected dummy pixel size per bucket. Output directly determines the safe \texttt{TARGET\_ALT} parameter.
    \item \textbf{Speed sweep} (\texttt{speed\_sweep.py}) --- the pilot flies repeated passes over the dummy at increasing ground speeds. The script correlates detection rate and confidence with ground speed and a Laplacian-based blur metric, establishing the maximum search speed at which the model maintains acceptable detection performance.
    \item \textbf{GPS accuracy} (\texttt{gps\_accuracy.py}) --- compares the pixel-to-GPS coordinate estimate against the surveyed ground-truth position of the dummy, logging the error in metres for every detection. Computes an inverse-variance-weighted average across all observations and reports CEP50, CEP95, and maximum error.
    \item \textbf{GPS drift} (\texttt{gps\_drift.py}) --- the drone hovers stationary above a known point. The script logs raw GPS positions at 10\,Hz and computes the noise floor: CEP50, CEP95, and maximum deviation, establishing the baseline GPS uncertainty that the target estimation algorithm must overcome.
    \item \textbf{Model comparison} (\texttt{model\_compare.py}) --- benchmarks all available \texttt{.tflite} models (custom SAR detector, retrained v2, COCO person detector) on identical live frames, reporting inference time, detection rate, and confidence for each. Informs model selection for the mission.
    \item \textbf{Detection log} (\texttt{detection\_log.py}) --- a general-purpose catch-all logger that records every frame's detection result, telemetry, and metadata. Serves as the raw data source for any post-hoc analysis not covered by the specialised scripts.
\end{enumerate}

% ─────────────────────────────────────────────────────────────────────
\subsection{Defect Discovery by Tier}\label{sec:bug-discovery}

Table~\ref{tab:bug-cost-appendix} lists representative defects, the tier at which they were discovered, the resolution time, and the estimated cost had they propagated to flight.

\begin{table}[H]
\centering
\caption{Defects discovered at each testing tier, with resolution time and estimated cost if undetected.}
\label{tab:bug-cost-appendix}
\small
\begin{tabular}{@{}p{4.5cm}clp{4.5cm}@{}}
\toprule
\textbf{Defect} & \textbf{Tier} & \textbf{Fix Time} & \textbf{Cost if Missed} \\
\midrule
Geofence sign convention inverted (\texttt{pointPolygonTest} positive = inside) & 1 & 15\,min & SSSI incursion; regulatory violation \\
\addlinespace
Repulsive force direction reversed at NFZ boundary & 1 & 20\,min & Drone pushed \emph{into} NFZ \\
\addlinespace
Barometric drift causes infinite landing loop & 1 & 30\,min & Battery depletion in hover \\
\addlinespace
Duplicate target queued from consecutive frames & 1 & 10\,min & Repeated descent on same target \\
\addlinespace
State transition missing timeout guard & 1 & 15\,min & Stuck in CENTERING indefinitely \\
\addlinespace
BGR/RGB colour inversion (IMX296) & 2 & 45\,min & Detection confidence drops to 0.3; mission failure \\
\addlinespace
FOV miscalibration (7.0 vs 5.46\,mm) & 2 & 1\,hr & GPS estimates 28\% off; landing outside 10\,m radius \\
\addlinespace
Python 3.13 pyserial breakage & 2 & 2\,hr & No Pi--Cube communication \\
\addlinespace
\texttt{tflite-runtime} unavailable on 3.13 & 2 & 30\,min & No inference on Pi \\
\addlinespace
Camera resource conflict (concurrent access) & 2 & 20\,min & Stream or detection crashes on startup \\
\addlinespace
GPS timing lag (100--200\,ms) & 3$^*$ & N/A & 0.8--1.6\,m systematic offset at 8\,m/s \\
\addlinespace
\texttt{self.investigating} uninitialised & 2 & 5\,min & Crash on first detection in \texttt{pi\_flight.py} \\
\bottomrule
\multicolumn{4}{@{}l}{\footnotesize $^*$Discovered via DJI video analysis, a proxy for Tier~3.} \\
\end{tabular}
\end{table}

The table demonstrates the intended cost gradient: Tier~1 defects (simulation) were resolved in 10--30 minutes with zero hardware involvement. Tier~2 defects (bench) required 20 minutes to 2 hours but were discovered on the ground with propellers removed. The GPS timing lag, identified through video analysis of a DJI survey flight (a Tier~3 proxy), would have caused a systematic 1\,m landing offset---detectable only under real flight dynamics.

% ─────────────────────────────────────────────────────────────────────
\subsection{V-Model Mapping}\label{sec:vmodel-mapping}

The five-tier framework maps onto the classical V-model~\cite{forsberg1991relationship} as follows:

\begin{table}[H]
\centering
\caption{Mapping between V-model verification levels and the progressive testing tiers.}
\label{tab:vmodel-map}
\small
\begin{tabular}{@{}lll@{}}
\toprule
\textbf{V-Model Level} & \textbf{Tier} & \textbf{Verification Activity} \\
\midrule
Requirements analysis & --- & Mission brief (R01--R12) + KML zones \\
System design & --- & State machine, architecture, module interfaces \\
Subsystem design & --- & \texttt{vision.py}, \texttt{planning.py}, \texttt{utils.py} APIs \\
Implementation & --- & Source code + unit tests \\
\midrule
Unit testing & 1 & Unit tests (\texttt{tests/unit/}, 6 scripts), TFLite model validation \\
Integration testing & 1, 2 & SITL end-to-end, bench command pipeline \\
System testing & 3, 4 & Passive flight data collection, log-only autonomous \\
Acceptance testing & 5 & Full autonomous mission with operator verification \\
\bottomrule
\end{tabular}
\end{table}

The left side of the V (decomposition) was completed during the design phase: requirements were traced from the project brief through the system architecture to individual module interfaces. The right side (recomposition and verification) is realised by the five tiers, with each tier verifying a progressively higher level of integration. The critical addition relative to a textbook V-model is the separation of system testing into \emph{passive} (Tier~3) and \emph{log-only autonomous} (Tier~4), which inserts two low-risk verification stages between integration testing and acceptance testing---stages that a strict V-model would collapse into a single system test phase.

% ─────────────────────────────────────────────────────────────────────
\subsection{Cost--Risk Gradient}\label{sec:cost-risk}

Table~\ref{tab:cost-gradient} quantifies the cost and risk at each tier, reinforcing the economic argument for catching defects early.

\begin{table}[H]
\centering
\caption{Cost and risk gradient across testing tiers.}
\label{tab:cost-gradient}
\small
\begin{tabular}{@{}clcccc@{}}
\toprule
\textbf{Tier} & \textbf{Environment} & \textbf{Hardware at Risk} & \textbf{Fix Turnaround} & \textbf{Iteration Cost} & \textbf{People Required} \\
\midrule
1 & Laptop (SITL) & None & Minutes & Electricity & 1 (developer) \\
2 & Bench (no props) & None & 30\,min--2\,hr & Travel to lab & 1--2 \\
3 & Passive flight & Airframe & Hours & Flight slot & 3+ (pilot, operator, spotter) \\
4 & Auto log-only & Airframe & Hours--day & Flight slot & 3+ \\
5 & Full autonomous & Airframe + target & Day+ & Flight slot + repair & 3+ \\
\bottomrule
\end{tabular}
\end{table}

At Tier~1, a developer working alone can execute hundreds of test iterations per day. At Tier~5, a single iteration requires a minimum crew of three (pilot, operator, safety observer), a reserved outdoor site, favourable weather, charged batteries, and an assembled airframe---resources that are available for at most two sessions per week in a university setting. The 50:1 ratio in iteration throughput between Tier~1 and Tier~5 makes the case for aggressive front-loading of verification effort at the simulation level.

\subsection{Summary}

The progressive methodology treated testing not as a phase that follows development, but as a continuous, tiered process that shapes development decisions. The five-tier gating system, supported by 58 dedicated scripts across eight categories (Table~\ref{tab:test-categories}), ensured that every integration boundary was verified before the system was permitted to cross it. Of the 15 significant defects discovered, 80\% were caught at Tiers~1 and~2 where the fix cost was measured in minutes, not field sessions. The structured experiment scripts prepared for Tier~3 passive flights provide a repeatable, quantitative basis for tuning configuration parameters before autonomy is enabled---replacing the ad-hoc ``fly and see'' approach with a data-driven verification pipeline.
